% Proved after full root and companion audits; exact actual-field candidate retained.
\subsection{Third-level point fields exclude the tame full-image comparison}

Throughout this subsection, assume for contradiction the remaining tame
comparison of B87--B88. Thus \(K\) is complete discretely valued of mixed
characteristic \((0,5)\), with algebraically closed residue field,
\(e_K=20\), and \(\mu_5\subset K\). Two whole potentially supersingular
height-two groups have isomorphic actual \(125\)-torsion, distinct finite
special inertia characters, and tame reducible residual inertia.
After the already justified common finite twist and interchange, take
\(B_1=25,B_2=75\), with high lines \(\mathbf F_5e_1\) and
\(\mathbf F_5e_2\). The common tame good field \(L/K\) has degree 24.
Put \(F=L(M[125])\), \(F_2=L(M[25])\). The actual comparison identifies
these fields for the two models. B87 gives
\[
 W=\operatorname{Gal}(F/L)=1+5M_2(\mathbf Z/125\mathbf Z),
 \qquad U=\operatorname{Gal}(F/F_2)=1+25M_2(\mathbf F_5).
 \tag{B-RD90.1}
\]
Every lower break is \(b(g)=v(g\varpi-\varpi)-1\), for the normalized
valuation on the top field of the indicated extension.

\begin{lemma}[The actual third-point towers and their two lower breaks]
\label{lem:B-RD90-point-towers}
\status{Proved; full independent root and companion audits passed.}
One can choose a common basis \(Q_1,Q_2\) of the actual \(125\)-torsion
whose images \(25Q_i=e_i\) lie on the two displayed residual axes and
whose point fields are stable under a tame lift. Put
\[
 P_i=5Q_i,\qquad E_i=L(P_i),\qquad E_i^+=L(Q_i).
\]
Then \(E_i^+/E_i\) is Galois of degree 25, and
\[
\begin{array}{c|c|c|c}
 &\text{lower breaks}&\text{diagonal cyclic quotient}&
       \text{off-diagonal cyclic quotient}\\\hline
 E_1^+/E_1&(606,6936)&1872&606\\
 E_2^+/E_2&(624,6834)&624&1866.
\end{array}                                                   \tag{B-RD90.2}
\]
The terms diagonal and off-diagonal refer to the actual translations
of \(Q_i\) by the first-level module, with respect to \(e_1,e_2\).
\end{lemma}
\begin{proof}
The finite extension \(F/K\) is Galois, with normal 5-group \(W\) and
cyclic quotient of order 24. A lift of a generator has order
\(24\cdot5^r\); a suitable power therefore gives a lift \(\tau\) of
order 24. Its action on the torsion lattice diagonalizes integrally:
the two distinct residual eigenvalues have unique Teichmuller lifts,
and the corresponding coprime factors of \(X^{24}-1\) give integral
idempotents. Choose \(Q_i\) on these two direct summands. Then
\(\tau Q_i\) is a unit multiple of \(Q_i\), so \(E_i,E_i^+\) are
stable under \(\tau\). This choice works in both models through the
actual torsion isomorphism.

Write a matrix in \(W\) as \(1+5X\), with
\(X\in M_2(\mathbf Z/25\mathbf Z)\). The stabilizers of \(P_i\)
and \(Q_i\) are respectively
\[
 H_i=\{1+5X:X e_i=0\pmod5\},\qquad
 S_i=\{1+5X:X e_i=0\pmod{25}\}.
 \tag{B-RD90.3}
\]
Their orders are \(5^6\) and \(5^4\). On \(H_i\), the map
\(g\mapsto gQ_i-Q_i\) takes values in the \(L\)-rational first-level
module. It is a homomorphism, and the terminal matrices \(U\) already
surject onto that module. Its kernel is \(S_i\). This proves that
\(E_i^+/E_i\) is Galois of degree 25, and that every nonzero translation
occurs. It does not assert that \(E_i^+/L\) is Galois.

The first-level low/high valuations over \(L\) are \((5,95)\) and
\((15,45)\). The exact inverse polygon (B-RD76.3), whose coefficient
bounds allow wild absolute index, applies twice. Already the first
inverse over a high target is uniform because \(B_i>e_K\); all later
inverses are uniform. Thus division by five twice divides each of
these valuations by \(25^2\). Since \([E_i^+:L]=625\), the coordinate
valuations in \(E_i^+\) equal the first-level valuations again.
Translation by a first-level point changes a formal coordinate by
that point's coordinate times an integral unit. The complete rows are
therefore
\[
\begin{array}{c|c|r|r|r|c}
 \text{field}&\text{model}&A&D_0&D_1&\mathcal H\\\hline
 E_1^+&1&95&3125&59375&\text{diagonal}\\
 E_1^+&2&15&9375&28125&\text{off-diagonal}\\
 E_2^+&1&5&3125&59375&\text{off-diagonal}\\
 E_2^+&2&45&9375&28125&\text{diagonal}.
\end{array}                                                   \tag{B-RD90.4}
\]
Here \(A=v_{E_i^+}(T(Q_i))\), while the displacement is \(D_1\) on
the line \(\mathcal H\) and \(D_0\) off it. All four values \(A\)
have 5-adic order one. The strict coefficient cutoff is
\(v_{E_i^+}(5)=625\cdot480=300000\), above every displacement in
the table.

Tame conjugation on these translation groups has the same two weights
as in B87: zero and 6 for \(i=1\), zero and 18 for \(i=2\), modulo
24. Indeed the eigenvalue of \(Q_i\) cancels in the cocycle
\(\tau g\tau^{-1}Q_i-Q_i\). On a cyclic degree-five quotient over
\(E_i\), the total ramification degree over \(L\) is 125, so its
leading tame uniformizer multiplier \(\eta\) satisfies
\(\eta^{125}=\zeta\). Thus \(\eta=\zeta^5\). The conjugation
calculation of B87 again gives a positive diagonal break divisible
by 24, and an off-diagonal break congruent to 6 or 18 respectively.
The two quotient breaks differ; the last ramification line is one
of the two eigenaxes. For lower breaks \(t<u\), their cyclic quotient
breaks are \(t\) and \((u+4t)/5\).

We give the coefficient exclusions explicitly. For any row of
(B-RD90.4), let \(J\) be the first nonzero coordinate exponent prime
to five. The proof of Lemma~\ref{lem:B-RD87-coefficient} gives the
first possible displacement degree
\[
                  \min\{A+5b,J+b\}.                 \tag{B-RD90.5}
\]
A strict minimum cannot cancel. If this minimum is strictly less
than the actual displacement, cancellation therefore requires
\(J=A+4b\); in every case \(J>A\) and \(5\nmid J\).
The value \(J=\infty\) is allowed before these constraints force
a finite exponent. All uses are strictly below the stated cutoff.

In \(E_1^+/E_1\), if the last line were off-diagonal, the first
row and Lemma~\ref{lem:B-RD87-coefficient} would give the larger
cyclic break \((3125-95)/5=606\). Hence \(t<606\) and
\(u=3030-4t\). In the second row,
\(15+5t<9375\), so cancellation forces \(J=15+4t\).
At \(u>t\) this prime term changes uniquely first, at
\(J+u=3045\), instead of 28125. This is impossible.
The last line is diagonal. The second row now gives its cyclic
break \((9375-15)/5=1872\), so
\[
             t\equiv6\pmod{24},\quad 0<t<1872,
             \qquad u=9360-4t.                      \tag{B-RD90.6}
\]
For the first row, if \(t<606\), then \(95+5t<3125\);
cancellation forces \(J=95+4t\), whose displacement at \(u\)
is 9455, instead of 59375. If \(t>606\), one must have
\(J=3125-t\). At \(u\), this term changes strictly before the
leading term, since
\[
 J-95-4u=15t-34410<0 \quad(t<1872).
\]
Its displacement is \(12485-5t<59375\), again impossible.
Only \(t=606,u=6936\) remains. The compulsory first prime
exponents in the two rows are then 27839 and 2439; these are
indeed prime to five. This proves the first row of (B-RD90.2).

In \(E_2^+/E_2\), a diagonal last line would give the larger
quotient break \((3125-5)/5=624\), so \(t<624\) and
\(u=3120-4t\). The fourth row forces \(J=45+4t\), whose
unique displacement at \(u\) is 3165, instead of 28125.
The last line must instead be off-diagonal. Its cyclic break
is \((9375-45)/5=1866\), so
\[
             t\equiv0\pmod{24},\quad 0<t<1866,
             \qquad u=9330-4t.                      \tag{B-RD90.7}
\]
For the third row, \(t<624\) forces \(J=5+4t\) and the
incorrect displacement 9335 at \(u\). If \(t>624\), it forces
\(J=3125-t\), and
\(J-5-4u=15t-34200<0\). The unique displacement is
\(12455-5t<59375\), again impossible. Thus \(t=624,u=6834\).
The two first prime exponents are 27341 and 2541. This proves
the second row of (B-RD90.2).
\end{proof}

\begin{proposition}[The forced terminal filtration in the actual full field]
\label{prop:B-RD90-terminal}
\status{Proved; full independent root and companion audits passed.}
Under the same contradictory hypotheses, the terminal extension
\(F/F_2\) would have lower breaks
\[
                7896,\quad11646,\quad769146,\quad937896,
 \tag{B-RD90.8}
\]
with the successive coordinate directions removed in the order
\(B_{22},B_{21},B_{12},B_{11}\), for \(1+25B\in U\).
For elements outside \(U\), the lower breaks in \(F/L\) would
equal their first-layer breaks in \(F_2/L\), namely
\(6,96,1146,1896\), in the coordinate order
\(A_{21},A_{22},A_{12},A_{11}\).
\end{proposition}
\begin{proof}
First compute the actual intersections, rather than imposing
independence of abstract extensions. The subgroups in (B-RD90.3)
satisfy \(US_i=H_i\). Consequently
\[
 F_2\cap E_i^+=E_i,\qquad D_i:=F_2E_i^+=F^{U\cap S_i}.
 \tag{B-RD90.9}
\]
The quotient \(H_i/(U\cap S_i)\) has order \(5^4\) and is
elementary abelian. For its first-layer matrices \(A,C\), both
\(Ae_i=Ce_i=0\), so \([A,C]e_i=0\). The exact commutator
and fifth-power identities modulo 125 are
\[
 [1+5A,1+5C]=1+25[A,C],\qquad
       (1+5A+25B)^5=1+25A.
 \tag{B-RD90.10}
\]
Both right sides lie in \(U\cap S_i\). Hence \(D_i/E_i\) is
the actual elementary-abelian compositum of the disjoint
degree-25 extensions \(F_2/E_i\) and \(E_i^+/E_i\).

The subgroup filtration of the already proved \(F_2/L\) profile
in B88 gives lower breaks \((96,1146)\) for \(F_2/E_1\), with
characters \(A_{22},A_{12}\), and \((6,1896)\) for \(F_2/E_2\),
with characters \(A_{21},A_{11}\). Their upper breaks are
\((96,306)\) and \((6,384)\), respectively. The new extensions
have upper breaks \((606,1872)\) and \((624,1866)\), by
(B-RD90.2). All four breaks in each compositum are distinct.
The unequal-break sum argument proved in B88 therefore gives
all its character conductors, without assuming any further
linear independence of coefficients. Exact upper-to-lower
integration yields
\[
\begin{array}{c|c|c}
 &\text{upper breaks of }D_i/E_i&\text{lower breaks of }D_i/E_i\\\hline
 i=1&(96,306,606,1872)&(96,1146,8646,166896)\\
 i=2&(6,384,624,1866)&(6,1896,7896,163146).
\end{array}                                                   \tag{B-RD90.11}
\]
The last two character directions in each row are precisely
those from \(E_i^+/E_i\). Restricting the lower filtration to
\(\operatorname{Gal}(D_i/F_2)\) and then passing to upper
numbering therefore gives
\[
\begin{array}{c|c|c|c}
 &\text{lower breaks}&\text{upper breaks}&\text{terminal characters}\\\hline
 D_1/F_2&(8646,166896)&(8646,40296)&B_{21},B_{11}\\
 D_2/F_2&(7896,163146)&(7896,38946)&B_{22},B_{12}.
\end{array}                                                   \tag{B-RD90.12}
\]
Here the characters are the restrictions of the actual
translation maps \(g\mapsto gQ_i-Q_i\) to \(U\), hence are
exactly the indicated matrix entries, not chosen abstract
identifications of isomorphic Galois groups.

Since \((U\cap S_1)\cap(U\cap S_2)=1\), one has
\(D_1D_2=F\); their degree-25 extensions over \(F_2\) are
linearly disjoint. Thus \(F/F_2\) is elementary abelian of
rank four with coordinate upper breaks
\[
       B_{22}:7896,\quad B_{21}:8646,\quad
       B_{12}:38946,\quad B_{11}:40296.
 \tag{B-RD90.13}
\]
Integrating indices \(1,5,25,125\) gives (B-RD90.8).
These are the actual lower subgroup breaks in \(F/L\),
because lower numbering is preserved on subgroups.

Every nonidentity element of \(U\) consequently has break
at least 7896. For any \(g\notin U\), the exact coset-average
formula for \(F_2/L\) shows that some lift of its image has
break less than 7896: otherwise the quotient break would
be at least 7896, contrary to B88. Multiplication by any
element of \(U\) preserves such a smaller break, by the
valuation inequality for composition of automorphisms.
All lifts therefore have the same break, equal to the
quotient break. This proves the final assertion and
completely identifies every nonidentity break in \(W\).

Only subgroup restriction and the exact norm-polynomial
quotient/Herbrand formulas are used here, with the full proofs
already read in \cite[Proposition 3.7, Theorem 3.10 and Proposition 3.13(ii), full proofs, pp.~14--16]
{ColmezCorpsLocaux}. No Hasse--Arf theorem or assertion about
unrelated extensions is required.
\end{proof}

\begin{theorem}[No different-special-character tame residual comparison at 125]
\label{thm:B-RD90-tame-exclusion}
\status{Proved; full independent root and companion audits passed.}
For \(p=5,e_K=20\) in the stated setting, two whole potentially
supersingular height-two groups with isomorphic actual
\(125\)-torsion and tame reducible residual inertia have the
same finite special inertia characters.
\end{theorem}
\begin{proof}
Otherwise the preceding reductions give the actual full image
(B-RD90.1) and its forced filtration. On the good field \(L\),
the determinant is the cyclotomic character. Its image is the
full group \(1+5\mathbf Z/125\mathbf Z\), of order 25.
The actual quotient field is therefore \(C=L(\mu_{125})\),
and the determinant kernel has order \(5^6=15625\).

Fix a nonidentity order-five element of this quotient, of
determinant \(1+25r\), \(r\in\mathbf F_5^\times\).
Write each element uniquely as \(g=1+5A+25B\), with entries
of \(A,B\) chosen from \(0,1,2,3,4\). A lift in this coset
must have \(\operatorname{tr}(A)=0\) modulo five. For each
such \(A\), exactly 125 matrices \(B\) have the specified
determinant: the second determinant coordinate is an affine
trace equation, including the carry in \(\operatorname{tr}(A)\)
and \(\det A\). No incorrect entrywise linearization of
the full determinant is used.

Among the 124 nonzero traceless first-layer matrices,
100 have \(A_{21}\ne0\), 20 have \(A_{21}=0,A_{22}\ne0\),
and four have only \(A_{12}\ne0\). Their breaks are
6, 96, and 1146. For \(A=0\), the equation is exactly
\(B_{11}+B_{22}=r\). The terminal filtration just proved
gives respectively 100, 20, four, and one solutions with
breaks 7896, 11646, 769146, and 937896. Hence the exact
coset counts are
\[
\begin{array}{c|rrrrrrr}
 b(g)&6&96&1146&7896&11646&769146&937896\\\hline
 \#g&12500&2500&500&100&20&4&1.
\end{array}                                                   \tag{B-RD90.14}
\]
Their sum is 15625 and their weighted break sum is 5925000.
For this nonidentity quotient element \(\bar g\), the exact
formula of \cite[Proposition 3.7 and full proof]{ColmezCorpsLocaux}
is
\[
 b_C(\bar g)+1
   =\frac{1}{15625}\sum_{g\mapsto\bar g}\bigl(b_F(g)+1\bigr).
\]
It would give
\[
                  b_C(\bar g)=\frac{5925000}{15625}
                              =\frac{1896}{5}.        \tag{B-RD90.15}
\]
This is impossible: \(b_C(\bar g)=v_C(\bar g\varpi_C-\varpi_C)-1\)
is an integer for the normalized discrete valuation of the
actual field \(C\). The contradiction proves the theorem.
\end{proof}

\begin{remark}[Scope and exact supplement]
\label{rem:B-RD90-scope}
\status{Proved; full independent root and companion audits passed.}
This excludes only the tame reducible residual branch left
by B87--B88. The irreducible-residual and nontrivial-residual-wild
branches retain their previously proved restrictions; no
actual comparison in those branches is constructed or excluded
here. In particular the unrestricted first-wild least-index
interval remains \(\{2,3\}\). The allowed base constructed
in B89 is unaffected: it never supplied the other division
fields or a height-two torsion comparison.

The retained exact program is
\path{code/reduction_descent/verify_p5_third_point_ramification.py}.
Its first run passed at 02:14:42.322512 UTC on 20 November 2026
in 0.08913616598874796 seconds. It checks all 206 lower-break
possibilities inside the four bounds proved above, obtaining
the two displayed survivors; it reconstructs every quotient,
subgroup, and compositum break by exact rational arithmetic.
It also enumerates all 62500 actual matrices in the four
nonidentity order-five determinant fibers, independently
evaluating their determinants modulo 125, and checks 1250
stabilizer commutators. These finite checks supplement the
field identifications and uniform coefficient argument;
they do not replace them or assert realization of a numerical
profile. No failed execution preceded the recorded run.
All code and output are retained. The complete B60, B76,
B87, and B88 arguments were reread; no new external-source
opening is claimed for this candidate.

\smallskip
\noindent\emph{Independent audits.}
The complete root and monodromy proof and exact-code audits passed.
Their literal author-output-preserving reruns took
\(0.08867679999093525\) and
\(0.08529232500586659\) seconds, respectively.
The root subsequently reread the original Colmez proofs through
Proposition 3.13 and retained a separate source addendum. The canonical
citation names Proposition 3.7, Theorem 3.10 and Proposition 3.13(ii)
precisely; this corrects only the candidate's bibliographic labels.
The exact candidate, source scopes, code, outputs and promotion diff
remain retained.
\end{remark}
